\documentclass[a4paper,12pt]{report}
\usepackage[utf8]{inputenc}
\usepackage[T1]{fontenc}
\usepackage[french]{babel}
\usepackage{geometry}
\geometry{margin=2.5cm}
\usepackage{graphicx}
\usepackage{enumitem}
\usepackage{array}
\usepackage{booktabs}
\usepackage{tabularx}
\usepackage{xcolor}
\usepackage{titlesec}
\usepackage{fancyhdr}
\usepackage{hyperref}
\usepackage{tikz}
\usepackage{tcolorbox}
\usepackage{longtable}

% Couleurs BPM
\definecolor{bpmred}{HTML}{C8102E}
\definecolor{bpmgreen}{HTML}{00703C}
\definecolor{bpmblue}{HTML}{1A3A5C}
\definecolor{bpmgray}{HTML}{F5F5F5}

% Titres
\titleformat{\chapter}[display]{\normalfont\Huge\bfseries\color{bpmred}}{\chaptertitlename\ \thechapter}{20pt}{\Huge}
\titleformat{\section}{\normalfont\Large\bfseries\color{bpmblue}}{\thesection}{1em}{}
\titleformat{\subsection}{\normalfont\large\bfseries\color{bpmblue!80}}{\thesubsection}{1em}{}

% En-tête/pied de page
\pagestyle{fancy}
\fancyhf{}
\fancyhead[L]{\small\textcolor{bpmblue}{Cahier des Charges — GAB Intelligence Dashboard}}
\fancyhead[R]{\small\textcolor{bpmblue}{Banque Populaire du Maroc}}
\fancyfoot[C]{\thepage}
\renewcommand{\headrulewidth}{0.4pt}

\hypersetup{
    colorlinks=true,
    linkcolor=bpmblue,
    urlcolor=bpmred,
}

\begin{document}

% ===== PAGE DE GARDE =====
\begin{titlepage}
\centering
\vspace*{1cm}

{\Large\textsc{Banque Populaire du Maroc}}\par
\vspace{0.5cm}
{\large Projet de Fin d'Études — 2024/2025}\par

\vspace{2cm}
\rule{\textwidth}{2pt}\par
\vspace{0.8cm}
{\Huge\bfseries\textcolor{bpmred}{Cahier des Charges}}\par
\vspace{0.5cm}
{\LARGE\textcolor{bpmblue}{Système de Détection et de Prédiction\\des Pannes des GAB}}\par
\vspace{0.3cm}
{\Large GAB Intelligence Dashboard}\par
\vspace{0.8cm}
\rule{\textwidth}{2pt}\par

\vspace{2cm}

\begin{tabular}{ll}
\textbf{Organisme d'accueil :} & Banque Populaire du Maroc \\
\textbf{Encadrant entreprise :} & \textit{[Nom de l'encadrant]} \\
\textbf{Encadrant académique :} & \textit{[Nom de l'encadrant]} \\
\textbf{Réalisé par :} & \textit{[Votre nom complet]} \\
\textbf{Filière :} & \textit{[Votre filière]} \\
\textbf{Année universitaire :} & 2024 — 2025 \\
\end{tabular}

\vfill
{\small Document confidentiel — Diffusion restreinte}
\end{titlepage}

% ===== TABLE DES MATIÈRES =====
\tableofcontents
\newpage

% ===================================================================
\chapter{Contexte Général du Projet}
% ===================================================================

\section{Présentation de l'organisme d'accueil}

La \textbf{Banque Populaire du Maroc (BPM)} est l'un des principaux groupes bancaires marocains, disposant d'un réseau étendu d'agences et de Guichets Automatiques Bancaires (GAB) couvrant l'ensemble du territoire national.

Le parc de GAB de la BPM compte environ \textbf{200 automates} répartis dans \textbf{13 villes} du Maroc, offrant des services de retrait, de consultation de solde et d'opérations bancaires 24h/24 et 7j/7.

\section{Problématique}

Les pannes des GAB représentent un enjeu majeur pour la banque~:

\begin{itemize}[leftmargin=2cm]
    \item \textbf{Impact financier} — Chaque panne engendre des coûts d'intervention corrective élevés (estimés à \textbf{5\,000 MAD} par intervention) et un manque à gagner sur les transactions.
    \item \textbf{Impact client} — L'indisponibilité d'un GAB dégrade l'expérience client et l'image de la banque.
    \item \textbf{Maintenance réactive} — Le modèle actuel repose sur une maintenance corrective (intervention après panne), générant des temps d'arrêt prolongés.
\end{itemize}

Le taux de panne moyen observé est de l'ordre de \textbf{9,8\%} sur l'historique disponible (2022–2023), avec des disparités géographiques et saisonnières significatives.

\section{Objectifs du projet}

Le projet vise à concevoir et développer un \textbf{système intelligent de détection et de prédiction des pannes} des GAB, permettant~:

\begin{enumerate}[leftmargin=2cm]
    \item \textbf{Prédire les pannes 48h à l'avance} grâce à des modèles de Machine Learning.
    \item \textbf{Visualiser en temps réel} l'état du parc GAB via un tableau de bord interactif.
    \item \textbf{Optimiser les coûts de maintenance} en passant d'un modèle correctif à un modèle préventif.
    \item \textbf{Fournir des outils d'aide à la décision} pour les équipes de maintenance et la direction.
\end{enumerate}

% ===================================================================
\chapter{Périmètre Fonctionnel}
% ===================================================================

\section{Modules fonctionnels}

Le système est structuré en \textbf{six modules} principaux~:

\begin{table}[h!]
\centering
\renewcommand{\arraystretch}{1.4}
\begin{tabularx}{\textwidth}{|c|l|X|}
\hline
\rowcolor{bpmblue!10}
\textbf{\#} & \textbf{Module} & \textbf{Description} \\
\hline
1 & Vue d'ensemble & KPIs globaux, tendances mensuelles, répartition par type/environnement/saison/âge \\
\hline
2 & Analyse géographique & Classement des villes par taux de panne, heatmap villes $\times$ mois \\
\hline
3 & Comparaison de modèles & Radar des métriques (F1, Précision, Rappel, AUC-ROC, AUC-PR), matrice de confusion \\
\hline
4 & Importance des features & Top features par importance, répartition par famille, export multi-format \\
\hline
5 & Optimisation du seuil & Simulation économique coûts/bénéfices, seuil optimal par F1 et par coût \\
\hline
6 & Scoring en temps réel & Prédiction de risque de panne par GAB, jauge interactive, contributions des variables \\
\hline
\end{tabularx}
\caption{Modules fonctionnels du système}
\end{table}

\section{Fonctionnalités transversales}

\begin{itemize}
    \item \textbf{Filtrage global} — Filtres multi-sélection par ville (13), type de GAB (Diebold, Hyosung, NCR, Wincor), et année.
    \item \textbf{Monitoring temps réel} — Surveillance continue avec alertes automatiques pour les niveaux ÉLEVÉ et CRITIQUE.
    \item \textbf{Export de données} — Export en CSV, Excel et PDF avec mise en forme professionnelle.
    \item \textbf{Thème clair/sombre} — Interface adaptable avec persistance du choix utilisateur.
    \item \textbf{Journal d'audit} — Traçabilité de toutes les prédictions effectuées (format JSONL).
\end{itemize}

\section{Niveaux de risque}

Le système classifie chaque GAB selon quatre niveaux~:

\begin{table}[h!]
\centering
\renewcommand{\arraystretch}{1.3}
\begin{tabular}{|l|c|l|l|}
\hline
\rowcolor{bpmblue!10}
\textbf{Niveau} & \textbf{Score} & \textbf{Couleur} & \textbf{Action recommandée} \\
\hline
FAIBLE & $< 0.15$ & \textcolor{green!60!black}{Vert} & Aucune action requise \\
\hline
MODÉRÉ & $0.15 - 0.35$ & \textcolor{blue}{Bleu} & Surveillance renforcée \\
\hline
ÉLEVÉ & $0.35 - 0.65$ & \textcolor{orange}{Orange} & Planifier une maintenance préventive \\
\hline
CRITIQUE & $\geq 0.65$ & \textcolor{red}{Rouge} & Intervention immédiate recommandée \\
\hline
\end{tabular}
\caption{Classification des niveaux de risque}
\end{table}

% ===================================================================
\chapter{Spécifications Techniques}
% ===================================================================

\section{Architecture du système}

Le système adopte une architecture \textbf{client-serveur} à deux tiers~:

\begin{tcolorbox}[colback=bpmgray, colframe=bpmblue, title=Architecture technique]
\begin{center}
\begin{tikzpicture}[
    box/.style={draw, rounded corners, minimum width=4cm, minimum height=1.2cm, align=center, font=\small},
    arrow/.style={->, thick, >=stealth}
]
\node[box, fill=blue!10] (frontend) at (0,0) {\textbf{Frontend}\\React 18 + Vite 5\\Port 5173};
\node[box, fill=red!10] (backend) at (7,0) {\textbf{Backend}\\FastAPI + Uvicorn\\Port 8000};
\node[box, fill=green!10] (data) at (7,-3) {\textbf{Données \& Modèles}\\CSV + Pickle + JSON};
\node[box, fill=orange!10] (ml) at (0,-3) {\textbf{Pipeline ML}\\scikit-learn\\5 modèles};

\draw[arrow] (frontend) -- node[above]{\footnotesize API REST} (backend);
\draw[arrow] (backend) -- (data);
\draw[arrow] (ml) -- (data);
\draw[arrow, dashed] (ml) -- (backend);
\end{tikzpicture}
\end{center}
\end{tcolorbox}

\section{Stack technologique}

\begin{table}[h!]
\centering
\renewcommand{\arraystretch}{1.3}
\begin{tabularx}{\textwidth}{|l|l|X|}
\hline
\rowcolor{bpmblue!10}
\textbf{Couche} & \textbf{Technologie} & \textbf{Rôle} \\
\hline
Frontend & React 18, Vite 5, Plotly.js & Interface utilisateur interactive, graphiques \\
\hline
Backend & FastAPI, Uvicorn & API REST, logique métier, scoring \\
\hline
Machine Learning & scikit-learn, Pandas, NumPy & Entraînement, évaluation, inférence \\
\hline
Sérialisation & Joblib (Pickle) & Persistance des modèles entraînés \\
\hline
Export & openpyxl, ReportLab, Matplotlib & Génération de rapports CSV/Excel/PDF \\
\hline
\end{tabularx}
\caption{Stack technologique}
\end{table}

\section{Données}

\subsection{Jeu de données}

\begin{itemize}
    \item \textbf{Volume} — 146\,000 observations sur 2 ans (janvier 2022 – décembre 2023).
    \item \textbf{Couverture} — 200 GAB répartis dans 13 villes marocaines.
    \item \textbf{Variable cible} — \texttt{panne\_sous\_48h} (binaire, taux positif $\approx$ 9,8\% — classe déséquilibrée).
    \item \textbf{Features} — 101 variables d'ingénierie réparties en 7 familles~:
\end{itemize}

\begin{table}[h!]
\centering
\renewcommand{\arraystretch}{1.2}
\begin{tabular}{|l|l|}
\hline
\rowcolor{bpmblue!10}
\textbf{Famille} & \textbf{Description} \\
\hline
Original & Variables brutes (âge, température, erreurs, transactions) \\
\hline
Lag (J-1/3/7) & Valeurs retardées à 1, 3 et 7 jours \\
\hline
Rolling (7-14j) & Moyennes et écarts-types glissants \\
\hline
Tendance & Dérivées et pentes temporelles \\
\hline
Temporel & Mois, jour de la semaine, saison \\
\hline
Catégoriel & Encodages (ville, type GAB, environnement) \\
\hline
Interaction & Croisements entre variables clés \\
\hline
\end{tabular}
\caption{Familles de features}
\end{table}

\subsection{Stratégie de découpage}

\begin{itemize}
    \item \textbf{Découpage temporel} — Entraînement sur 2022, test sur 2023 (pas de fuite temporelle).
    \item \textbf{Calibration} — Régression isotonique pour calibrer les probabilités.
\end{itemize}

\section{Modèles de Machine Learning}

Cinq modèles sont entraînés et comparés~:

\begin{enumerate}
    \item Régression Logistique
    \item Random Forest
    \item Gradient Boosting
    \item Support Vector Machine (SVM)
    \item Réseau de Neurones (MLP)
\end{enumerate}

Un modèle \textit{Dummy} (stratifié) sert de référence de base. Le meilleur modèle est automatiquement sélectionné selon le score F1.

\section{Mode dégradé (Heuristique)}

En cas d'indisponibilité du modèle ML, un \textbf{algorithme heuristique} prend le relais avec 9 composantes pondérées dérivées de l'importance des features~:

\begin{table}[h!]
\centering
\renewcommand{\arraystretch}{1.2}
\begin{tabular}{|l|c|}
\hline
\rowcolor{bpmblue!10}
\textbf{Composante} & \textbf{Poids} \\
\hline
Jours depuis dernière maintenance & 22\% \\
Score de négligence maintenance & 14\% \\
Ratio d'erreurs & 14\% \\
Tendance température & 11\% \\
Température absolue & 9\% \\
Âge du GAB & 9\% \\
Risque matériel & 7\% \\
Taux d'erreurs transactions & 6\% \\
Tendance erreurs & 3\% \\
\hline
\end{tabular}
\caption{Pondérations du mode heuristique}
\end{table}

\section{API REST — Endpoints}

\begin{longtable}{|l|l|p{6cm}|}
\hline
\rowcolor{bpmblue!10}
\textbf{Méthode} & \textbf{Endpoint} & \textbf{Description} \\
\hline
\endfirsthead
\hline
\rowcolor{bpmblue!10}
\textbf{Méthode} & \textbf{Endpoint} & \textbf{Description} \\
\hline
\endhead
GET & \texttt{/api/health} & État du système et mode du modèle \\
\hline
GET & \texttt{/api/metadata} & Métadonnées (villes, types, années) \\
\hline
GET & \texttt{/api/overview} & KPIs et tendances (filtrable) \\
\hline
GET & \texttt{/api/geography} & Analyse géographique et heatmap \\
\hline
GET & \texttt{/api/models} & Métriques de tous les modèles \\
\hline
GET & \texttt{/api/features} & Importance des variables \\
\hline
GET & \texttt{/api/threshold} & Simulation de seuil économique \\
\hline
POST & \texttt{/api/scoring} & Scoring en temps réel d'un GAB \\
\hline
GET & \texttt{/api/monitoring} & État temps réel de tous les GAB \\
\hline
GET & \texttt{/api/predictions-stats} & Statistiques du journal d'audit \\
\hline
POST & \texttt{/api/reload-model} & Rechargement à chaud du modèle \\
\hline
GET & \texttt{/api/export/*} & Export CSV / Excel / PDF \\
\hline
\caption{Endpoints de l'API REST}
\end{longtable}

% ===================================================================
\chapter{Exigences Non Fonctionnelles}
% ===================================================================

\begin{itemize}
    \item \textbf{Performance} — Temps de réponse API $<$ 500~ms pour les endpoints principaux. Cache avec TTL configurable (défaut 300s, monitoring 25s).
    \item \textbf{Disponibilité} — Mode dégradé heuristique garantissant le scoring même sans modèle ML.
    \item \textbf{Sécurité} — Token d'administration pour les opérations sensibles (rechargement de modèle). Variables d'environnement pour la configuration en production.
    \item \textbf{Traçabilité} — Journal JSONL de toutes les prédictions avec horodatage.
    \item \textbf{Ergonomie} — Interface responsive, thème clair/sombre, graphiques interactifs.
    \item \textbf{Maintenabilité} — Architecture modulaire, séparation front/back, code documenté.
\end{itemize}

% ===================================================================
\chapter{Livrables}
% ===================================================================

\begin{table}[h!]
\centering
\renewcommand{\arraystretch}{1.4}
\begin{tabularx}{\textwidth}{|c|X|l|}
\hline
\rowcolor{bpmblue!10}
\textbf{\#} & \textbf{Livrable} & \textbf{Format} \\
\hline
1 & Code source du tableau de bord (frontend + backend) & Répertoire Git \\
\hline
2 & Modèles ML entraînés et calibrés & Fichiers .pkl \\
\hline
3 & Jeu de données avec features d'ingénierie & CSV \\
\hline
4 & Documentation technique (API, architecture) & PDF \\
\hline
5 & Rapport de PFE & PDF \\
\hline
6 & Script de déploiement automatisé & Batch (.bat) \\
\hline
\end{tabularx}
\caption{Liste des livrables}
\end{table}

% ===================================================================
\chapter{Planning Prévisionnel}
% ===================================================================

\begin{table}[h!]
\centering
\renewcommand{\arraystretch}{1.3}
\begin{tabularx}{\textwidth}{|c|X|c|}
\hline
\rowcolor{bpmblue!10}
\textbf{Phase} & \textbf{Description} & \textbf{Durée} \\
\hline
1 & Étude de l'existant et analyse des besoins & 2 semaines \\
\hline
2 & Collecte et préparation des données & 2 semaines \\
\hline
3 & Ingénierie des features et entraînement ML & 3 semaines \\
\hline
4 & Développement du backend (API REST) & 3 semaines \\
\hline
5 & Développement du frontend (Dashboard) & 3 semaines \\
\hline
6 & Intégration, tests et optimisation & 2 semaines \\
\hline
7 & Rédaction du rapport et soutenance & 2 semaines \\
\hline
\end{tabularx}
\caption{Planning prévisionnel}
\end{table}

% ===================================================================
\chapter{Contraintes et Risques}
% ===================================================================

\section{Contraintes}

\begin{itemize}
    \item Données historiques limitées à 2 ans (2022–2023).
    \item Classe cible déséquilibrée ($\approx$ 9,8\% de pannes).
    \item Environnement de développement local (pas de serveur de production dédié).
\end{itemize}

\section{Risques identifiés}

\begin{table}[h!]
\centering
\renewcommand{\arraystretch}{1.3}
\begin{tabularx}{\textwidth}{|X|c|X|}
\hline
\rowcolor{bpmblue!10}
\textbf{Risque} & \textbf{Probabilité} & \textbf{Mitigation} \\
\hline
Qualité insuffisante des données & Moyenne & Nettoyage rigoureux, validation croisée \\
\hline
Surapprentissage des modèles & Moyenne & Découpage temporel, régularisation \\
\hline
Dérive du modèle en production & Élevée & Monitoring des prédictions, rechargement à chaud \\
\hline
Indisponibilité du modèle ML & Faible & Mode dégradé heuristique automatique \\
\hline
\end{tabularx}
\caption{Matrice des risques}
\end{table}

\end{document}
